\documentclass[12pt]{article}
\usepackage[T1]{fontenc}
\usepackage[utf8]{inputenc}
\usepackage[brazil]{babel}
\usepackage[margin=1in]{geometry}
\usepackage{ragged2e}
\usepackage[normalem]{ulem}
\setlength{\parindent}{0pt}
\setlength{\parskip}{0.8\baselineskip}
\begin{document}
\justifying
\noindent Verifica-se que a questão objeto da controvérsia jurídica suscitada no recurso manejado pela parte recorrente esteve afetada no representativo da controvérsia - \textbf{Tema n. }\textbf{230} da TNU - (PU no processo n. PEDILEF 0028697-44.2016.4.01.3900/PA),  afetado em 18/09/2019 foi julgado (trânsito em julgado em 01/06/2021), por meio do qual foi elucidada a seguinte questão:

\noindent \textit{"Estabelecer qual a base de cálculo do Imposto sobre Operações de Crédito, Câmbio e Seguro ou relativas a Títulos ou Valores Mobiliários (IOF) no caso de contratos de crédito prorrogados, renovados ou renegociados."}

\noindent Restou firmada a seguinte tese:

\noindent \textit{"Não haverá incidência de IOF complementar sobre o saldo devedor não liquidado de operação de crédito objeto de prorrogação, renovação, novação, composição, consolidação, confissão de dívida e negócios assemelhados, sem substituição de devedor, caso na operação de origem tenha sido aplicado o limite máximo previsto no art. 7º, §1º do Decreto n. 6306/2007 (alíquota vigente aplicada ao valor do principal colocado a disposição do devedor, multiplicada por 365 dias, acrescida da alíquota adicional de 0,038\%). Todavia, nos casos em que na operação de origem a alíquota aplicada tenha sido inferior à máxima prevista no Decreto n. 6.306/2007 haverá a incidência da exação, sobre o saldo não liquidado, sem que se cogite bis in idem. Por sua vez, a base de cálculo do IOF nos casos de contratos de crédito prorrogados, renovados ou renegociados é o saldo não liquidado. A entrega ou colocação de novos valores ao mutuário na mesma oportunidade constitui nova base de cálculo que permite a incidência de IOF nos termos do art. 7o § 9o do 6.306 de 14 de dezembro de 2007." (Tema 230 TNU).}

\noindent Com efeito, depreende-se da leitura do voto condutor do acórdão que o mesmo se encontra em conformidade com o decidido no  tema pela TNU.

\noindent A proposito, consta do acordao hostilizado: (...) Cédula de crédito bancário, Espécies de títulos de crédito, Obrigações, DIREITO CIVIL

\noindent Nessas condições, o conteúdo do Acórdão recorrido é consentâneo com o entendimento da TNU, firmado em julgamento de recurso representativo de controvérsia, o que atrai a incidência da regra prevista pelo art. 14, III, \textit{b,} da Resolução CJF n. 586/2019:

\noindent \textit{\textbf{Art. 14.}}\textit{ Decorrido o prazo para contrarrazões, os autos serão conclusos ao magistrado responsável pelo exame preliminar de admissibilidade, que deverá, de forma sucessiva: (...) }\textit{\textbf{III }}\textit{- negar seguimento a pedido de uniformização de interpretação de lei federal interposto contra acórdão que esteja em conformidade com entendimento consolidado: }\textit{\textbf{b)}}\textit{ em }\uline{\textit{\textbf{recurso representativo de controvérsia pela Turma Nacional de Uniformização}}}\textit{ ou em pedido de uniformização de interpretação de lei dirigido ao Superior Tribunal de Justiça;(...).}

\noindent Posto isso, com arrimo no \uline{\textbf{art. 14, III, }}\uline{\textit{\textbf{b}}}\textit{,} da Resolução CJF n. 586/2019 c/c art. 1.030, I, do CPC, \textbf{NEGO SEGUIMENTO }\textbf{a}o\textbf{ PU}\textbf{.}

\noindent Intimem-se. Aguarde-se o decurso do prazo recursal. Após, certifique-se o trânsito em julgado e devolva-se ao juízo de origem.
\end{document}
